% ID: SEL-D0043-OA-E02
% ESTADO: validada. S044, cotejada con examen y criterios.
\textbf{a)} Los tramos son derivables en sus interiores.
En $2$, la continuidad exige $4+2a=2+b$.
Las derivadas son $2x+a$ y $-(1+b)/(x-1)^2$,
así que las pendientes laterales coinciden si $4+a=-1-b$.
Sustituyendo $b=2+2a$ en esta última igualdad,
resulta $3a=-7$. Por tanto $\boxed{a=-7/3,\ b=-8/3}$.
Se cumplen continuidad y coincidencia de pendientes.

\textbf{b)} En $1$ se usa el primer tramo. Con $a=2$,
$f(1)=3$ y $f'(1)=4$. La tangente es $y-3=4(x-1)$,
es decir, $\boxed{y=4x-1}$.
